\section{評価実験}\label{sec5}

本評価実験では，家庭環境知識を把握していなければ解決することが困難である家庭内タスクのデータセットを用いて，提案手法によるタスクの成功率を評価する．

\subsection{評価用データセット}\label{sec5.1}
本研究では提案手法を評価するために，\ref{sec3}で示した要件を満たす新たなデータセットを構築した．具体的には，物体の状態に着目した「状態変化タスク」と，物体間の関係に着目した「配置タスク」の2種類のタスクを構築した．

\subsubsection{状態変化タスク}
本タスクでは``Turn on all lights''や``Close all windows''などの指示に対して，環境内の複数の物体を特定の状態に変更することが求められる．対象物には，スイッチのある家電や扉のある家具など，状態を有する物体が含まれる．

タスクの難易度は，一つの物体のみを扱う単純なタスク（例：``Turn on the light''）と，複数の物体を扱う複雑なタスク（例：``Turn on all lights''）の二つのカテゴリに分類される．

本研究では，後者のより複雑なタスクの設計に注力し，各タスク記述に対して，初期状態の異なる複数のパターンでデータセットを作成した．例えば``Turn on all lights''というタスクでは，初期状態がすべてオン，すべてオフ，またはオンとオフが混在している場合など，様々なバリエーションが存在する．エージェントは異なる初期状態に適応して行動する必要があるため，どの物体とインタラクションすべきかを判断するために，詳細な家庭環境知識の理解が求められ，タスクの難易度が高まる．

\subsubsection{配置タスク}
本タスクでは``Put all apples in the fridge''などの指示に対して，複数の物体を環境内の特定の場所に配置することが求められる．対象物には，ロボットが把持可能な物体と，他の物体を受け入れることができる物体（``receptacle''）が含まれる．

状態変化タスクと同様に，一つの物体を配置する単純なタスク（例：``Put an apple in the fridge''）と，複数の物体を配置する複雑なタスク（例：``Put all apples in the fridge''）の二つのカテゴリに分類される．

ここでも，本研究では複雑なタスクの設計に焦点を当てた．例えば冷蔵庫のような``receptacle''には，``OPEN''または``CLOSED''という扉の状態がある．データセットには，``receptacle''の初期状態として``OPEN''と``CLOSED''の2種類のパターンが含まれる．そのため，物体を配置する前に，エージェントはその``receptacle''の扉を開ける必要があるかどうかを家庭環境知識に基づいて推論する必要がある．このような設計によって，環境に対する推論なしにはタスクを完遂できず，タスクの難易度が高くなる．

\subsubsection{データセット構築}
\begin{figure}[tb]
    \begin{lstlisting}[caption=状態変化タスクのデータセット例, label=state-change-task]
{
    "task": "Turn on all lightswitches",
    "scene": 1,
    "initial_room": "kitchen",
    "initial_states": [
        {"id": 71, "states": ["ON"]},
        {"id": 173, "states": ["OFF"]},
        {"id": 261, "states": ["ON"]},
        {"id": 427, "states": ["OFF"]}
    ],
    "goal_states": [
        {"id": 71, "states": ["ON"]},
        {"id": 173, "states": ["ON"]},
        {"id": 261, "states": ["ON"]},
        {"id": 427, "states": ["ON"]}
    ],
    "action_scripts": [
        "[WALK] <lightswitch> (173)",
        "[SWITCHON] <lightswitch> (173)",
        "[WALK] <lightswitch> (427)",
        "[SWITCHON] <lightswitch> (427)"
    ]
}
    \end{lstlisting}
\end{figure}

\begin{figure}[tb]
    \begin{lstlisting}[caption=配置タスクのデータセット例, label=placement-task]
{
    "task": "Put all plums in the fridge",
    "scene": 4,
    "initial_room": "bedroom",
    "initial_states": [
        {"id": 103, "states": ["CLOSED"]}
    ],
    "goal_states": [
        {
            "from_id": 53,
            "to_id": 103,
            "relation_type": "INSIDE"
        },
        {
            "from_id": 54,
            "to_id": 103,
            "relation_type": "INSIDE"
        }
    ],
    "action_scripts": [
        "[WALK] <plum> (53)",
        "[GRAB] <plum> (53)",
        "[WALK] <fridge> (103)",
        "[OPEN] <fridge> (103)",
        "[PUTIN] <plum> (53) <fridge> (103)",
        "[WALK] <plum> (54)",
        "[GRAB] <plum> (54)",
        "[WALK] <fridge> (103)",
        "[PUTIN] <plum> (54) <fridge> (103)",
        "[CLOSE] <fridge> (103)"
    ]
}
    \end{lstlisting}
\end{figure}

本研究では，\ref{sec4.6}で述べたように，シミュレータとしてVHを使用する．したがって，データセットはVHとの互換性を持つように設計した．

VHには7種類の異なるシーン（家屋）があり，それぞれに複数の部屋とそこに設置された様々な物体が含まれている．各タスクは，これらシーン毎の特性に応じて構成している．

VHにおいて，物体の状態は``ON''，``OFF''，``OPEN''，``CLOSED''の4種類が存在する．VH内の物体のバリエーションを考慮し，状態変化タスクのためのデータセットは，電源を持つ7種類の物体（例：``lightswitch''，``tablelamp''）を対象として構築した．さらに，配置タスクにおいては，飲食物（例：``banana''，``milk''）といったカテゴリに属する物体のみを対象とし，それらの配置場所（``receptacle''）も，それに適した場所（例：``kitchentable''，``fridge''）に限定した．このデータセットには，12種類の飲食物と8種類の``receptacle''が含まれる．

データセットの例を，Listing~\ref{state-change-task}および~\ref{placement-task}に示す．データセットは，タスクの記述，シーン，エージェントの初期位置，物体の初期状態，物体の目標状態，そしてタスクを完遂するために必要なアクションスクリプトで構成される．

\begin{figure*}[tb]
\footnotesize
\begin{align}
 \label{formula5.1}
 Success Rate(SR) &= \frac{Sum Of Successes}{Sum Of Tasks}\\
 \label{formula5.2}
 Action Execution Failure Rate(AEFR) &= \frac{Sum Of Action Execution Failures}{Sum Of Tasks}\\
 \label{formula5.3}
 Failure Rate Of Reaching Maximun Attempts(FRRMA) &= \frac{Sum Of Reaching Maximun Attempts}{Sum Of Tasks}\\
 \label{formula5.4}
 Erroneous Terminate Failure Rate(ETFR) &= \frac{Sum Of Erroneous Terminate Failures}{Sum Of Tasks}\\
 \label{formula5.5}
 Failure Rate(FR) &= AEFR + FRRMA + ETFR\\
 \label{formula5.6}
 Average Steps &= \frac{Sum Of Action Steps For All Tasks}{Sum Of Tasks}
\end{align}
\end{figure*}


\subsubsection{データの分割方法}
データの分割方法について説明をする．状態変化タスクデータセットの事例数は464個，配置タスクデータセットの事例数は154個であり，テスト用と訓練用データセットの事例数が大体2対1の割合になるように，各データセットを分割する．

訓練用データセットは，~\ref{sec4.3}で述べた，プロンプトに与えるタスク例に対するアクションスクリプトの出力事例として用いる．
状態変化タスクのテスト用事例数は312個，訓練用事例数は152個，配置タスクのテスト用事例数は103個，訓練用事例数は51個とした．

これらのデータセットは，同名タスクでもオブジェクトの初期状態や，VHのシーンを変更して，データによってタスクを遂行するために必要な振る舞いが異なるように作成をしているため，別のタスクとして扱う．そのため，分割はランダムに行ったが，同名タスクのデータが，テスト用と訓練用で混在しないように分割をしたため，正確に2対1の割合での分割はできていない．

\subsection{評価方法}\label{sec5.2}
実験の前段階として，データセットのシーンと初期状態を用いてVH環境の初期設定をする．次に提案手法にデータセットのタスク記述を入力して，タスクを実行するための適切な行動計画を生成する．

提案手法の実験では，アクションステップを生成するときのLLMへのプロンプトの与え方を二つの方法で行う．一つ目は，\ref{sec4.3}で説明したプロンプトの複数の内容を１回で与える方法である．二つ目は，\ref{sec4.3}で説明したプロンプトの複数の内容を，内容ごとに分割して複数回に分けてLLMに与える方法である．さらに，提案手法の構成要素の一つである「前提条件の検証」の効果を検証するために，アブレーションスタディを実施する．他の構成要素については，それを除くと提案手法が実行不可能となるためアブレーションスタディは実施していない．

アクションステップの生成時に用いる，タスクを遂行するためのアクションスクリプトの生成例の選択方法について説明する．例の選択は，訓練データセットの中から，入力となるタスク名と類似度の高いタスク名を1種類選択する．タスク名に対して複数個のデータが存在する場合には，タスクを遂行するためのアクションスクリプトの生成例が最長のものを選択する．

次に\ref{sec4.4}で述べた試行回数の上限の決定方法を説明する．試行回数の上限は，データセットに含まれるタスクを遂行するためのアクションスクリプトの例の長さの2倍と設定し，設定した上限までにタスクを達成できなければ，強制終了して，タスク失敗とする．また，提案手法では，LLMがタスクを達成したと判断し，``End''を出力するまで，アクションの実行を継続する．試行回数が上限に達して，強制終了した場合にゴール状態をすべて満たしている場合も考えられるが，タスクを達成していると判断できず，アクションを実行し続けようとしたとして，タスク失敗とする．

また，ベースライン手法として，家庭環境知識を用いないHuangらの手法~\cite{huang22}を評価する．これは，本研究において構築したデータセットが，家庭環境知識を用いないと解くことが困難であることを実証するためのものである．この手法は，オブジェクトのIDの特定はできないため，自動でオブジェクトを探索するシミュレータの設定を採用する．

本評価実験では，GPT-4o mini (gpt-4o-mini-2024-07-18)とGPT-4o (gpt-4o-2024-08-06) の\HL{2種類のLLMを用いて}実験を行う．なお，ベースライン手法で用いるLLMは，GPT-4oのみとしている．GPT-4oはGPT-4o miniと比較して高性能なモデルであり，その出力をベースラインとすることで，提案手法との比較を妥当に行うことができる．このことから，ベースラインとしてはGPT-4oのみで十分であると判断した．

\subsection{評価指標}\label{sec5.3}

本研究では，既存研究\cite{ALFRED20,song2023llm}の評価で用いられているタスクの成功率によって提案手法を評価する．
これらの研究での評価では，物体の識別番号を考慮せず評価を行っているが，本評価では，物体の識別番号の一致も考慮して評価を行う．

成功率は，式\ref{formula5.1}に示すように，タスクの総数に対する成功したタスクの割合として定義する．タスクの成功条件は，ゴール状態のすべてを満たした場合に限り成功とし，一部でも満たしていない場合は失敗とする．

失敗には三つのタイプがある．一つ目はシミュレータにおけるアクション実行による失敗，二つ目は設定された最大試行回数に達したことによる失敗，三つ目はタスク完了と判定したが，タスクのゴール状態をすべて満たせていないことによる失敗である．失敗の原因を分析するために，各失敗タイプの割合を式~\ref{formula5.2}から式~\ref{formula5.4}に定義する．

式\ref{formula5.2}は，シミュレータにおけるアクション実行の失敗に起因する失敗率を示しており，タスクの総数に対するアクションの実行に失敗した回数の割合として定義する．式\ref{formula5.3}は，設定された最大試行回数に達したことによる失敗率を示しており，タスクの総数に対する当該失敗回数の割合として定義する．式\ref{formula5.4}は，タスク完了と判定されたにも関わらず，タスクのゴール状態を完全には満たしていないことによる失敗率を示しており，タスクの総数に対する当該失敗回数の割合として定義する．

また，式~\ref{formula5.6}により``$AverageSteps$''（アクションステップ数の平均値）を求める．ここでアクションステップ数は，アクションをシミュレータで実行した回数である．


\subsection{評価結果}\label{sec5.4}
\begin{table*}[tb]
    \caption{「状態変化タスク」の評価結果  ※PreCheck:「前提条件の検証」}
    \label{result1}
    \centering
    \begin{tabular}{llccccc}
    \hline
         \multirow{2}{*}{Method} & \multirow{2}{*}{LLM} & \multirow{2}{*}{SR} & \multicolumn{3}{c}{FR} & \multirow{2}{*}{Average Steps}\\
          & & & AEFR & FRRMA & ETFR\\
    \hline
         Baseline~\cite{huang22} & GPT-4o & 0.032 & 0.000 & 0.535 & 0.433 & 10.481\\
    \hline    
         \multirow{2}{*}{Single Prompt} & GPT-4o mini & 0.625 & 0.026 & 0.231 & 0.119 & 4.747\\
         & GPT-4o & 0.788 & 0.016 & 0.173 & 0.022 & 4.381\\
         
         \multirow{2}{*}{Multi Prompts} & GPT-4o mini & 0.750 & 0.045 & 0.183 & 0.022 & 4.917\\
         & GPT-4o & \textbf{0.894} & 0.019 & 0.000 & 0.087 & 3.391\\
    \hline
        \multirow{2}{*}{Single Prompt w/o PreCheck} & GPT-4o mini & 0.615 & 0.003 & 0.253 & 0.128 & 4.872\\
         & GPT-4o & 0.760 & 0.016 & 0.189 & 0.035 & 4.449\\
         
         \multirow{2}{*}{Multi Prompts w/o PreCheck} & GPT-4o mini & 0.696 & 0.016 & 0.253 & 0.035 & 5.080\\
         & GPT-4o & 0.872 & 0.010 & 0.016 & 0.103 & 3.478\\
    \hline
    \end{tabular}
\end{table*}

\begin{table*}[tb]
    \caption{「配置タスク」の評価結果  ※PreCheck:「前提条件の検証」}
    \label{result2}
    \centering
    \begin{tabular}{llccccc}
    \hline
         \multirow{2}{*}{Method} & \multirow{2}{*}{LLM} & \multirow{2}{*}{SR} & \multicolumn{3}{c}{FR} & \multirow{2}{*}{Average Steps}\\
          & & & AEFR & FRRMA & ETFR\\
    \hline
         Baseline~\cite{huang22} & GPT-4o & 0.000 & 0.010 & 0.573 & 0.417 & 15.398\\
    \hline    
         \multirow{2}{*}{Single Prompt} & GPT-4o mini & 0.466 & 0.146 & 0.029 & 0.359 & 7.282\\
         & GPT-4o & 0.738 & 0.117 & 0.029 & 0.117 & 8.078\\
         
         \multirow{2}{*}{Multi Prompts} & GPT-4o mini & 0.408 & 0.117 & 0.029 & 0.447 & 6.447\\
         & GPT-4o & \textbf{0.816} & 0.117 & 0.019 & 0.049 & 8.272\\
    \hline
        \multirow{2}{*}{Single Prompt w/o PreCheck} & GPT-4o mini & 0.476 & 0.117 & 0.019 & 0.388 & 6.592\\
         & GPT-4o & 0.680 & 0.165 & 0.029 & 0.126 & 7.320\\
         
         \multirow{2}{*}{Multi Prompts w/o PreCheck} & GPT-4o mini & 0.437 & 0.126 & 0.010 & 0.427 & 6.136\\
         & GPT-4o & 0.796 & 0.117 & 0.029 & 0.058 & 7.767\\
    \hline
    \end{tabular}
\end{table*}

提案手法とベースライン手法の比較結果について，状態変化タスクデータセットに基づく結果を\ref{result1}に，配置タスクデータセットに基づく結果を\ref{result2}に示す．

\ref{result1}と\ref{result2}における``Method''列の``Baseline''はHuangらの手法~\cite{huang22}を指し，``Single Prompt''は単一のプロンプトをLLMに提供するアプローチ，``Multi Prompts''はプロンプトを複数回に分けてLLMに提供するアプローチを指す．``Single Prompt w/o PreCheck''および``Multi Prompt w/o PreCheck''は，提案手法から「前提条件の検証」を除去したアブレーションスタディの結果である．

これらの結果から，提案手法はベースライン手法よりも高い成功率を達成することが示された．これは，ベースライン手法が家庭環境知識を考慮せずにアクションスクリプトを生成するため，家庭内タスクの完了が難しいことを示唆している．さらに，提案手法の``Single Prompt''と``Multi Prompts''を比較すると，プロンプトを複数のステップに分割することで，一般的に成功率が向上することが分かった．

アブレーションスタディの結果からは，「前提条件の検証」を適用することで，成功率が一般的に向上することが示された．しかし，GPT-4o miniを配置タスクに使用した場合には逆の効果が観察された．

また，\ref{result1}と\ref{result2}を比較すると，配置タスクの成功率は状態変化タスクよりも全体的に低いことが分かる．さらに，配置タスクにおけるGPT-4o miniとGPT-4oの結果の顕著な違いから，このタスクはより高度な推論能力を要求することが分かった．

\subsection{考察}\label{sec5.5}
状態変化タスクと配置タスクを比較すると，配置タスクの成功率が低いことが分かる．この違いは，LLMが単にオブジェクトの状態を識別するよりも，オブジェクト間の空間関係を認識し，推論することの難しさが影響している可能性がある．特に，配置タスクにおけるGPT-4o miniとGPT-4oの顕著な性能差は，このタスクがより高度な推論能力を必要とすることを示唆している．

状態変化タスクでは，全体の失敗率（FR）の中で試行回数が上限に達したことによる失敗の割合（FRRMA）が高いことが確認された．この問題は，家庭環境知識に基づく「アクションステップ生成」のプロセスに関連している可能性が高い．特に，目標状態をすでに満たしている物体や操作が不要な物体に対して不要なアクションを生成してしまうことで，提案手法が過剰なアクションステップを実行し，最終的に失敗するケースが多い．そのため，試行回数を増やさずとも失敗を予測する技術開発が今後の課題となる．

配置タスクでは，全体のFRの中でタスクの完了判定の誤判定による失敗の割合（ETFR）が高い．このことから，タスク完了の正確な判断が重要であることを示している．これらの失敗は，家庭環境知識とタスク記述に基づいてタスク完了を判断する「タスク完了の判定」の失敗に起因している．そのため，家庭環境知識とタスク記述に基づいてタスク完了を判断する「タスク完了の判定」の精度向上が今後の課題となる．なお，この考察は，GPT-4o miniを用いた配置タスクにおいて，``Multi Prompts''が``Single Prompt''を上回らなかった理由の説明にもなる．

状態変化タスクに関するアブレーションスタディでは，「前提条件の検証」を適用することで成功率が向上した．これは主にFRRMAの減少によるものである．例えば，電源を入れるタスクにおいて，電源がすでにオンになっている場合，提案手法が不要な「電源を入れる」アクションを生成してしまう可能性がある．しかし，前提条件の検証を適用することで，このような冗長なアクションを検出し防ぐことができ，より効率的なタスク実行につながる．ただし，VH環境では，すでにオンになっている電源を再度オンにするアクションを実行してもエラーが発生しないため，アクション実行による失敗の割合（AEFR）には，ほとんど影響を与えなかった．

一方，配置タスクにおいて前提条件チェックの影響は，GPT-4oではより顕著に現れ，特に``Single Prompt''手法においてAEFRの減少に寄与した．しかし，GPT-4o miniでは顕著な効果は見られなかった．これは，前提条件チェックが適用される前の段階である「アクションステップ生成」の時点で失敗が発生しているためである．


\subsection{限界} \label{sec5.6}
本研究の三つの限界について述べる．

一つ目に，本研究で構築したデータセットにはいくつかの制限がある．本提案手法を評価するために2種類のタスクを導入しているが，特定のタスクタイプに制約されている．また物体の状態と関係の範囲も限られている．より汎用的な評価をするためには，より幅広いタスクを含むようにデータセットを拡張することが不可欠である．

二つ目の限界は，家庭環境の知識がどのように獲得されるかに関するものである．本研究では，家庭環境の知識はシミュレータの内部データから正確に得られると仮定している．しかし，実世界を対象とするアプリケーションでは，センサーやカメラ等を用いた環境認識能力が，基本的な要件となる．

最後に，アクションの実行に関する制限がある．現在のところ，すべての行動はシミュレートされた環境内で実行されるが，LLMの出力がロボットの行動を導く実世界のシナリオでは，誤動作や制御のミス（例：物体をつかみ損ねる）の可能性があり，重大な動作不良や事故つながる可能性がある．このため，LLMをロボット工学に安全に応用するためには，さらなる課題が生じる．